\documentclass[11pt,letterpaper]{article}

\usepackage[top=1in, bottom=1in, left=1in, right=1in, headheight=14pt]{geometry}
\usepackage{fontspec}
\setmainfont{DejaVu Serif}
\setsansfont{DejaVu Sans}
\setmonofont{DejaVu Sans Mono}

\usepackage{xcolor}
\usepackage{titlesec}
\usepackage{fancyhdr}
\usepackage{enumitem}
\usepackage{hyperref}
\usepackage{booktabs}
\usepackage{tabularx}
\usepackage{longtable}
\usepackage{colortbl}
\usepackage{multirow}
\usepackage{array}
\usepackage{parskip}
\usepackage{tcolorbox}
\usepackage{graphicx}
\usepackage{lastpage}

%% COLOR SCHEME — 80s Pop History: Deep Violet / Electric Indigo / Synth Rose
\definecolor{primary}{HTML}{4A148C}
\definecolor{secondary}{HTML}{1565C0}
\definecolor{tertiary}{HTML}{880E4F}
\definecolor{accent}{HTML}{7B1FA2}
\definecolor{lightprimary}{HTML}{F3E5F5}
\definecolor{lightsecondary}{HTML}{E3F2FD}
\definecolor{lighttertiary}{HTML}{FCE4EC}
\definecolor{rowgray}{HTML}{F5F5F5}
\definecolor{bordergray}{HTML}{BDBDBD}
\definecolor{subtlegray}{HTML}{757575}
\definecolor{darktext}{HTML}{212121}

%% SECTION FORMATTING
\titleformat{\section}
  {\Large\bfseries\sffamily\color{primary}}
  {\thesection}{1em}{}
  [\vspace{-0.5em}\textcolor{bordergray}{\rule{\textwidth}{0.4pt}}]

\titleformat{\subsection}
  {\large\bfseries\sffamily\color{secondary}}
  {\thesubsection}{1em}{}

\titleformat{\subsubsection}
  {\normalsize\bfseries\sffamily\color{tertiary}}
  {\thesubsubsection}{1em}{}

\titlespacing*{\section}{0pt}{1.5em}{0.8em}
\titlespacing*{\subsection}{0pt}{1.2em}{0.5em}
\titlespacing*{\subsubsection}{0pt}{0.8em}{0.3em}

%% CUSTOM ENVIRONMENTS
\newcommand{\stageheader}[2]{%
  \noindent\colorbox{#2}{\makebox[\textwidth][l]{%
    \textcolor{white}{\bfseries\sffamily\hspace{6pt}#1\hspace{6pt}}}}%
  \vspace{0.5em}\par%
}

\tcbuselibrary{skins,breakable}

\newtcolorbox{asciibox}{
  colback=rowgray, colframe=bordergray,
  arc=1pt, boxrule=0.5pt,
  left=6pt, right=6pt, top=4pt, bottom=4pt,
  fontupper=\small\ttfamily,
  breakable
}

\newtcolorbox{quotebox}{
  colback=lightprimary, colframe=primary,
  arc=2pt, boxrule=0.5pt,
  left=12pt, right=12pt, top=8pt, bottom=8pt,
  breakable
}

\newcommand{\metafield}[2]{\textbf{\sffamily #1} #2\par}

%% TABLE HELPERS
\newcolumntype{L}[1]{>{\raggedright\arraybackslash}p{#1}}
\newcolumntype{C}[1]{>{\centering\arraybackslash}p{#1}}
\newcommand{\tableheader}[1]{\textbf{\sffamily\textcolor{white}{#1}}}

%% HEADERS / FOOTERS
\pagestyle{fancy}
\fancyhf{}
\fancyhead[L]{\small\sffamily\textcolor{subtlegray}{Duran Duran}}
\fancyhead[R]{\small\sffamily\textcolor{subtlegray}{GSD Mission Package}}
\fancyfoot[C]{\small\sffamily\textcolor{subtlegray}{Page \thepage\ of \pageref{LastPage}}}
\renewcommand{\headrulewidth}{0.4pt}
\renewcommand{\footrulewidth}{0pt}

%% HYPERLINKS
\hypersetup{
  colorlinks=true,
  linkcolor=secondary,
  urlcolor=secondary,
  pdfauthor={GSD Ecosystem},
  pdftitle={Duran Duran -- Mission Package},
  pdfsubject={Vision to Mission Pipeline},
}

%% ============================================================
\begin{document}

%% TITLE PAGE
\thispagestyle{empty}
\begin{center}
\vspace*{3cm}

{\small\sffamily\textcolor{primary}{\textbf{GSD RESEARCH MISSION PACKAGE}}}

\vspace{0.3cm}
\textcolor{bordergray}{\rule{8cm}{0.4pt}}

\vspace{1.5cm}
{\fontsize{28}{34}\selectfont\bfseries\sffamily\textcolor{primary}{Duran Duran\\[0.3em]Sound, Vision, and the Architecture of Pop}}

\vspace{0.8cm}
{\Large\sffamily\textcolor{secondary}{Birmingham 1978 to the Rock and Roll Hall of Fame}}

\vspace{0.5cm}
{\large\sffamily\itshape\textcolor{tertiary}{A Deep Research Study}}

\vspace{3cm}
{\sffamily\textcolor{accent}{Vision $\rightarrow$ Research $\rightarrow$ Mission}}\\[0.3em]
{\small\sffamily\textcolor{accent}{Full Pipeline Execution}}

\vspace{3cm}
{\small\sffamily\textcolor{subtlegray}{Date: March 26, 2026 \quad|\quad Status: Ready for GSD Orchestrator}}\\[0.3em]
{\small\sffamily\textcolor{subtlegray}{Activation Profile: Squadron (12 roles) \quad|\quad Estimated: 6--8 context windows across 3 sessions}}
\end{center}
\newpage

\tableofcontents
\newpage

%% ============================================================
\stageheader{STAGE 1 --- VISION DOCUMENT}{primary}

\section{Duran Duran --- Vision Guide}

\metafield{Date:}{March 26, 2026}
\metafield{Status:}{Research Complete / Mission Ready}
\metafield{Depends On:}{New Romantic movement history, MTV launch context, Second British Invasion scholarship}
\metafield{Context:}{Cold-start deep research mission. No prior conversation. Full three-stage pipeline.}

\subsection{Vision}

There is a moment, somewhere in 1982, when a new medium and a new band collide in a way that reshapes how music is experienced. Duran Duran did not invent the music video. They did not invent the synthesizer, the art-school pop band, or the notion that a rock performance could be a fashion statement. What they did was \textit{architect} a system --- a complete audio-visual-cultural package --- at precisely the moment that system could propagate globally via a young cable channel called MTV.

The band that emerged from Birmingham's Rum Runner nightclub was built on a principle that this research mission will explore in full: \textit{the space between the notes matters as much as the notes themselves}. Duran Duran understood that the gap between sound and image, between music and fashion, between British post-punk austerity and the exotic glamour of locations like Sri Lanka and Antigua, was not empty space. It was where meaning lived. Russell Mulcahy shot their videos as mini-films. Patrick Nagel painted their album covers as high art. Malcolm Garrett designed their typography as visual architecture. And beneath it all, a rhythm section of John Taylor and Roger Taylor drove dance floors on two continents.

This mission documents the full system: how it was assembled, why it worked, what it produced, and what it left behind. Over 100 million records sold. Top-five UK albums in five consecutive decades. A 2022 Rock and Roll Hall of Fame induction. And a generation of bands --- The Killers, Franz Ferdinand, The Bravery --- who have cited Duran Duran as the architectural blueprint for their own sound.

The Amiga Principle applies here perfectly: specialized execution paths, not brute force. Five members who each embodied a different discipline --- singer, keyboardist, bassist, guitarist, drummer --- and two managers who understood that a nightclub was a laboratory. Small, principled, elegant. Staggering in its yield.

\subsection{Problem Statement}

\begin{enumerate}
  \item \textbf{The Visual Revolution Problem.} Prior to 1981, music videos were largely promotional afterthoughts --- low-budget performance clips. Duran Duran, working with director Russell Mulcahy, transformed the format into cinematic mini-films shot on 35mm with professional crews, exotic locations, and deliberate narrative storytelling. This fundamentally altered the economics and aesthetics of the music industry, but the full scope of their contribution to visual music has never been systematically documented.

  \item \textbf{The Transatlantic Positioning Problem.} EMI UK promoted Duran Duran as New Romantic, a genre barely known in the US. Capitol Records (EMI's American subsidiary) had no framework for selling them. The band's eventual US breakthrough --- requiring a strategic remixing of Rio as a ``dance album,'' plus MTV's heavy rotation of their visually stunning videos --- represents a fascinating case study in international cultural positioning that deserves full analysis.

  \item \textbf{The Longevity Paradox.} Bands built primarily on teen idol appeal (the ``Fab Five,'' Princess Diana's favorites, covers of every teen magazine) typically burn out within a decade. Duran Duran achieved top-five UK albums in five consecutive decades. Understanding how they outlasted their contemporaries requires examining the substance beneath the style.

  \item \textbf{The Lineup Instability Question.} The band has never officially disbanded, yet underwent massive personnel changes --- the departures of Andy and Roger Taylor in 1986, the decade-long Warren Cuccurullo era, the 2001 reunion, Andy Taylor's cancer revelation and absence from the 2022 Rock and Roll Hall of Fame induction. How does a band maintain identity across such turbulence?

  \item \textbf{The Critical Undervaluation Problem.} As Moby noted, Duran Duran suffered the ``Bee Gees curse'': write amazing songs, sell 100 million records, and incur critical dismissal because of it. The gap between their commercial record and their critical reputation represents an unresolved tension in pop music scholarship that this mission addresses directly.
\end{enumerate}

\subsection{Core Concept}

\textbf{One-Line Model:} Duran Duran = (Post-Punk Architecture) $\times$ (New Romantic Aesthetics) $\times$ (MTV Distribution) $\times$ (Fashion-Forward Identity) --- the first complete audio-visual pop system.

This model captures the key insight: none of these four elements alone explains the band's success. Contemporaries had some combination of them. The Cure had post-punk architecture but rejected the fashion-forward mainstream. Spandau Ballet had the New Romantic aesthetics but never cracked America at the same scale. The Human League had the synthesizer sound but lacked the cinematic video ambition. Duran Duran assembled all four and deployed them simultaneously.

\subsubsection{Birmingham Club Scene Map}

\begin{asciibox}
BIRMINGHAM CLUB ECOSYSTEM (1978-1980)
======================================

BARBARELLA'S NIGHTCLUB
  |-- Named the band (Dr. Durand Durand from Barbarella, 1968)
  |-- Early gigs: Sex Pistols and The Clash played here
  |-- Nick Rhodes and John Taylor: regulars, forming ideas

RUM RUNNER NIGHTCLUB (Paul & Michael Berrow, owners)
  |-- Nick Rhodes: DJ (GBP 10/night)
  |-- John Taylor: Doorman
  |-- Roger Taylor: Busboy
  |-- Rehearsal space provided free
  |-- Paul & Michael Berrow became managers (Tritec Music, June 1980)
  |-- Classic lineup first show: July 9, 1980

HOLE IN THE WALL PUB (now Saramoons)
  |-- John Taylor + Nick Rhodes discuss band name over drinks, 1978

AIR STUDIOS LONDON (George Martin's studio)
  |-- Debut album recorded, Dec 1980 - Jan 1981
  |-- Rio recorded, early 1982

INFLUENCE STACK:
  David Bowie -> Roxy Music -> Kraftwerk -> Chic (Nile Rodgers)
  Japan -> Ultravox -> Human League -> Blondie
\end{asciibox}

\subsubsection{Module Architecture}

\begin{asciibox}
RESEARCH MODULE ARCHITECTURE
======================================

MODULE A: ORIGINS & FORMATION (1978-1981)
  Birmingham roots / New Romantic context / Lineup evolution
  Rum Runner lab / EMI signing / Debut album

MODULE B: MUSICAL ARCHITECTURE (1981-1993)
  Sound design: synths + guitars + funk bass
  Album analysis: Duran Duran -> Rio -> SATR -> Wedding Album
  Production evolution / Night Versions / remixing philosophy

MODULE C: VISUAL REVOLUTION (1981-1985)
  Russell Mulcahy collaboration / 35mm film innovation
  Sri Lanka trilogy / Patrick Nagel / Malcolm Garrett
  MTV strategy / video as art form

MODULE D: CULTURAL IMPACT & LEGACY (1982-2026)
  Second British Invasion / Beatlemania-scale fame
  Lineup turbulence / comebacks / Rock and Roll Hall of Fame
  Influence on The Killers, Franz Ferdinand, et al.

CROSS-MODULE CONNECTIONS:
  A -> B: Rum Runner club scene -> dance/synth sound philosophy
  A -> C: Art school origins -> visual priority
  B -> C: Night Versions -> MTV rotation strategy
  B -> D: Wedding Album comeback -> critical reappraisal
  C -> D: Video as art form -> lasting visual legacy
\end{asciibox}

\subsection{Research Modules}

\subsubsection{Module A: Origins and Formation (1978--1981)}
Birmingham, England: the unlikely cradle of New Romantic glamour. This module covers the band's formation at the intersection of post-punk's collapse and the emergence of synthesizer-driven pop, the Rum Runner nightclub as creative laboratory and management incubator, the complex lineup evolution from Stephen Duffy through to Simon Le Bon, and the signing with EMI following a bidding war with Phonogram. Includes the debut album's production, reception, and the cultural context of the New Romantic scene.

\subsubsection{Module B: Musical Architecture (1981--1993)}
The sound of Duran Duran is a system: John Taylor's funk-inspired bass lines, Andy Taylor's rock-influenced guitar work (AC/DC and Van Halen through a New Wave filter), Nick Rhodes's synthesizer arrangements drawing on Kraftwerk and Eno, Roger Taylor's driving rhythm work, and Simon Le Bon's melodic vocal approach with cryptic, imagistic lyrics. This module analyzes the full discography through the Wedding Album (1993), tracing the production philosophy, the Night Versions remixing strategy for dance clubs, and the Nile Rodgers collaboration on The Reflex.

\subsubsection{Module C: Visual Revolution (1981--1985)}
Duran Duran were among the first bands to shoot music videos on 35mm film with professional movie crews, treating each clip as a ``mini-film'' rather than a performance document. Director Russell Mulcahy (who directed the first video ever played on MTV: The Buggles' ``Video Killed the Radio Star'') directed 10+ Duran Duran videos, including the three-part Sri Lanka trilogy for the Rio album. Album designer Malcolm Garrett applied Swiss Style modernism and postmodern pastiche to their sleeve art. Artist Patrick Nagel's illustration for the Rio cover became one of the most iconic images in pop history. This module documents the full visual system.

\subsubsection{Module D: Cultural Impact and Legacy (1982--2026)}
From the Second British Invasion of 1982--83 (when 30\% of US record sales were British acts) to the 2022 Rock and Roll Hall of Fame induction, this module traces the full arc of Duran Duran's cultural significance. Covers the Beatlemania-scale mania of 1983--84, the decline and reinvention through the late 1980s, the Wedding Album comeback, the 2001 reunion, and the modern legacy. Documents influence on bands including The Killers, Franz Ferdinand, and The Bravery. Includes Andy Taylor's prostate cancer announcement at the Hall of Fame ceremony.

\subsection{Chipset Configuration}

\begin{asciibox}
chipset:
  name: DuranDuran-Research-v1
  activation_profile: squadron
  modules:
    - id: MOD_A
      name: origins_formation
      model: sonnet
      priority: foundation
    - id: MOD_B
      name: musical_architecture
      model: sonnet
      priority: core
    - id: MOD_C
      name: visual_revolution
      model: opus
      priority: core
    - id: MOD_D
      name: cultural_legacy
      model: opus
      priority: synthesis
  synthesis:
    model: opus
    inputs: [MOD_A, MOD_B, MOD_C, MOD_D]
    output: complete_research_reference
  token_budget:
    opus: 30%
    sonnet: 60%
    haiku: 10%
\end{asciibox}

\subsection{Scope Boundaries}

\subsubsection{In Scope (v1.0)}
\begin{itemize}
  \item Band history from 1978 Birmingham formation through 2026
  \item Full discography analysis: 16 studio albums, key singles and extended versions
  \item Music video production history and aesthetic analysis
  \item New Romantic movement context and Second British Invasion
  \item All five classic lineup members' roles and contributions
  \item Post-1986 lineups including Warren Cuccurullo era
  \item Awards, chart records, and commercial data
  \item Influence on subsequent artists and critical reappraisal
  \item Rock and Roll Hall of Fame 2022 induction
  \item Cultural context: MTV launch, UK club scene, Birmingham music ecosystem
\end{itemize}

\subsubsection{Out of Scope}
\begin{itemize}
  \item Solo discographies of individual members in full depth (Arcadia, Power Station, Simon Le Bon solo)
  \item Bootleg recordings or unauthorized releases
  \item Detailed financial accounts or business litigation
  \item Personal biographies beyond band-relevant context
  \item Comprehensive tour setlist documentation
  \item Scholarly musicological score analysis
\end{itemize}

\subsection{Success Criteria}

\begin{enumerate}
  \item Origins module documents all lineup phases (1978--1980) with sources, including the Duffy, Wickett, and Thomas vocalist periods.
  \item The Rum Runner nightclub's role as both creative incubator and management origin is fully traced with specifics.
  \item Full discography covers all 16 studio albums with chart positions, key singles, and production notes.
  \item Sound architecture analysis identifies the specific musical roles and influences of each classic-lineup member.
  \item Visual Revolution module documents the Russell Mulcahy collaboration with specific production details for at least 5 key videos.
  \item Patrick Nagel's Rio cover and Malcolm Garrett's design system are analyzed as unified visual identity.
  \item MTV strategy documented with specific chart data showing video-rotation-to-sales correlation.
  \item Second British Invasion context provides quantitative data (chart positions, market share percentages).
  \item Legacy module cites at least 3 named contemporary bands acknowledging Duran Duran's influence.
  \item All chart claims attributed to specific sources (Billboard, UK Singles Chart official data).
  \item Lineup turbulence documented with dates, reasons, and musical consequences for all major personnel changes.
  \item Andy Taylor's 2022 Rock and Roll Hall of Fame absence and cancer announcement documented.
\end{enumerate}

\subsection{Relationship to Other Documents}

\begin{center}
\rowcolors{2}{rowgray}{white}
\begin{tabularx}{\textwidth}{L{4cm}XC{2.5cm}}
\toprule
\rowcolor{primary}
\tableheader{Document} & \tableheader{Relationship} & \tableheader{Status} \\
\midrule
New Romantic Movement Study & Contextual parent: ideological and aesthetic origins & Upstream \\
Second British Invasion Analysis & Overlapping scope: Duran Duran as case study within larger phenomenon & Sibling \\
MTV Cultural History & Critical context: the distribution system that made video strategy viable & Upstream \\
80s Pop Production Techniques & Technical sibling: production methods of the era & Sibling \\
GSD Amiga Vision & Philosophical connection: specialized execution paths = Amiga Principle & Downstream \\
\bottomrule
\end{tabularx}
\end{center}

\subsection{The Through-Line}

The Amiga Principle holds that architectural elegance beats brute force --- that specialized execution paths, faithfully iterated, produce staggering complexity from small, principled building blocks. Duran Duran is one of pop music's purest examples of this principle applied at the cultural scale.

They did not have the biggest budget, the best voice, the most technical musicianship, or the deepest critical credibility. What they had was a complete system: each member a specialized node, each collaborator (Mulcahy, Nagel, Garrett, Berrow brothers) occupying a distinct architectural role, and a distribution medium (MTV) that their visual-first philosophy was perfectly architected to exploit. The space between the sound and the image, between Birmingham and Sri Lanka, between post-punk rawness and New Romantic glamour --- that space was not empty. It was the system. It was where the meaning lived.

Forty-five years later, with 100 million records sold and a Hall of Fame plaque, Duran Duran remains proof that the most enduring outcomes are not built by the loudest or the fastest. They are built by the most architecturally coherent.

\newpage

%% ============================================================
\stageheader{STAGE 2 --- RESEARCH REFERENCE}{secondary}

\section{Duran Duran --- Research Reference}

\subsection{How to Use This Document}

This reference compiles verified findings from primary and secondary sources for all four research modules. Agents should use this as the evidentiary foundation for the final research document. All claims require source attribution; any unsourced claim encountered during execution should be flagged and verified before inclusion.

\textbf{Key source organizations and publications:}
\begin{itemize}
  \item Billboard (chart data and commercial records)
  \item Official UK Singles Chart / OCC (UK chart positions)
  \item Rolling Stone magazine (critical context, Second British Invasion analysis)
  \item RIAA (US certification data)
  \item Rock and Roll Hall of Fame (institutional recognition)
  \item Wikipedia (synthesis of secondary sources, cross-referenced)
  \item Classic Pop Magazine (album analysis, making-of features)
  \item The Saturday Evening Post (long-form retrospective)
  \item DuranDuran.com (official timeline and primary band sources)
  \item AllMusic / Rate Your Music (critical reception data)
\end{itemize}

\subsection{Module A Reference: Origins and Formation (1978--1981)}

\subsubsection{Founding Context}

John Taylor (born Nigel Taylor, 1960, Birmingham) and Nick Rhodes (born Nicholas Bates, 1962, Birmingham) formed Duran Duran in 1978, naming the band after the villain ``Dr. Durand Durand'' from Roger Vadim's 1968 science fiction film \textit{Barbarella}, the day after the film was broadcast on BBC (October 20, 1978). The name was conceived during a conversation over drinks at the Hole in the Wall Pub (now Saramoons) in Birmingham.

Initial members included Stephen Duffy (vocals and bass) and Simon Colley (bass and clarinet). Their first live performance was April 5, 1979, at Birmingham Polytechnic. Duffy and Colley departed June 1979.

\subsubsection{Lineup Evolution (1978--1980)}

\begin{center}
\rowcolors{2}{rowgray}{white}
\begin{tabularx}{\textwidth}{L{2.5cm}L{5cm}XC{2cm}}
\toprule
\rowcolor{primary}
\tableheader{Year} & \tableheader{Members} & \tableheader{Notable Event} & \tableheader{Source} \\
\midrule
1978 & Duffy, J. Taylor, Colley, Rhodes & First rehearsals; Hole in the Wall pub & DuranDuran.com \\
Apr 1979 & Duffy, J. Taylor, Colley, Rhodes & First live show, Birmingham Poly & DuranDuran.com \\
Jun 1979 & Andy Wickett (vocals), R. Taylor joins & Wickett co-wrote early ``Girls on Film'' & Wikipedia \\
Late 1979 & Jeff Thomas replaces Wickett & Brief tenure; Curtis also departs & Wikipedia \\
Apr 1980 & Andy Taylor joins from Newcastle & Melody Maker ad: ``Live-Wire Guitarist'' & DuranDuran.com \\
May 1980 & Simon Le Bon joins & Recommended by Rum Runner barmaid & DuranDuran.com \\
Jul 1980 & Classic ``Fab Five'' lineup set & First show: July 9, Rum Runner & DuranDuran.com \\
\bottomrule
\end{tabularx}
\end{center}

\subsubsection{The Rum Runner Laboratory}

The Rum Runner nightclub (owners: Paul and Michael Berrow) functioned simultaneously as performance venue, rehearsal space, employment source, and management incubator. John Taylor worked the door; Nick Rhodes DJed; Roger Taylor worked as a busboy. The Berrows formed Tritec Music (named after a triangular-themed bar within the club) in June 1980, signed the band to a production deal, and became their managers --- creating an unusually integrated early infrastructure that gave the band stability and continuity rare for a newly-formed act.

Simon Le Bon auditioned in May 1980, arriving in pink leopard-skin trousers with a notebook of lyrics. His lyric notebook contained what would become ``The Chauffeur.'' The band wrote their first song together, ``Sound of Thunder,'' at the Rum Runner immediately after Le Bon joined.

\subsubsection{EMI Signing and Debut Album}

After attracting attention supporting Hazel O'Connor on tour (November--December 1980), Duran Duran were the subject of a bidding war between EMI and Phonogram. They signed with EMI in December 1980. Their debut album was recorded in six weeks at AIR Studios (George Martin's London facility), produced by Colin Thurston, who had previously worked with David Bowie, Iggy Pop, the Human League, and Magazine.

The debut album, \textit{Duran Duran}, was released June 15, 1981. It reached \#3 on the UK Albums Chart and remained in the UK Top 100 for 117--118 weeks (sources vary slightly), achieving platinum status by December 1982. The initial US release failed to generate interest until after \textit{Rio}'s success triggered a reissue in 1983, which reached \#10 on the Billboard 200.

The album's sound, as writer Annie Zaleski describes it, featured ``space-age keyboards, post-punk guitars, disco-inspired bass lines and Le Bon's vocal croon'' --- drawing on Bowie, Roxy Music, Ultravox, Japan, Giorgio Moroder, Chic, and Kraftwerk.

\subsection{Module B Reference: Musical Architecture (1981--1993)}

\subsubsection{The Sound System: Member Roles}

\begin{center}
\rowcolors{2}{rowgray}{white}
\begin{tabularx}{\textwidth}{L{3cm}L{3.5cm}XL{2.5cm}}
\toprule
\rowcolor{primary}
\tableheader{Member} & \tableheader{Role} & \tableheader{Influences / Approach} & \tableheader{Sonic Contribution} \\
\midrule
Simon Le Bon & Lead vocals & Drama student; imagery, ambiguity & Melodic croon; cryptic lyrics \\
Nick Rhodes & Keyboards, synths & Kraftwerk, Bowie, Eno, Moroder & Arpeggiator synths; atmosphere \\
John Taylor & Bass guitar & Chic, Nile Rodgers, funk & Driving funk bass lines \\
Andy Taylor & Lead guitar & AC/DC, Van Halen, Jeff Beck & Rock guitar through NW filter \\
Roger Taylor & Drums, percussion & Steady, dance-floor driven & Foundation for club appeal \\
\bottomrule
\end{tabularx}
\end{center}

Nick Rhodes: ``He would do Eno and I would do Jeff Beck or Jimmy Page'' (Andy Taylor, describing his and Rhodes's contrasting approaches). John Taylor believed Andy Taylor's rock influences were key to their American breakthrough --- providing the familiar guitar energy that US audiences expected, filtered through a New Wave aesthetic.

\subsubsection{Key Studio Albums: Chart and Production Data}

\begin{center}
\rowcolors{2}{rowgray}{white}
\begin{tabularx}{\textwidth}{L{4cm}C{1.5cm}C{1.5cm}C{1.5cm}X}
\toprule
\rowcolor{primary}
\tableheader{Album} & \tableheader{Year} & \tableheader{UK} & \tableheader{US} & \tableheader{Key Notes} \\
\midrule
\textit{Duran Duran} & 1981 & \#3 & \#10 & 117+ weeks UK Top 100; platinum UK \\
\textit{Rio} & 1982 & \#2 & \#6 & 129 weeks US chart; double platinum US \\
\textit{Seven and the Ragged Tiger} & 1983 & \#1 & \#8 & Only UK \#1 album; ``The Reflex'' (US \#1) \\
\textit{Notorious} & 1986 & \#16 & \#12 & First without full classic lineup \\
\textit{Big Thing} & 1988 & \#15 & \#24 & Electronic shift; moderate success \\
\textit{Liberty} & 1990 & \#8 & \#46 & Reunion of J. Taylor; transition period \\
\textit{Duran Duran} (Wedding Album) & 1993 & \#4 & \#7 & ``Ordinary World'' (UK \#6, US \#3); comeback \\
\textit{Astronaut} & 2004 & \#3 & \#36 & Classic reunion; ``(Reach Up for The) Sunrise'' \\
\textit{All You Need Is Now} & 2011 & \#6 & \#10 & Critical re-engagement; Mark Ronson prod. \\
\textit{Future Past} & 2021 & \#5 & \#18 & Top five UK, fifth consecutive decade \\
\textit{Danse Macabre} & 2023 & --- & --- & Halloween-themed; latest release \\
\bottomrule
\end{tabularx}
\end{center}

\textit{Sources: Billboard, OCC / UK Albums Chart official data, AllMusic}

\subsubsection{The Night Versions Strategy}

Before digital sampling made remixing standard practice, Duran Duran were pioneers of the ``Night Version'' --- entirely re-recorded, extended performances of their singles, specifically designed for nightclub play and available on 12'' vinyl. Roger Taylor recalled: ``We always did the Night Versions of the songs for clubs, but at that time we had to record them live.'' These versions were compiled on \textit{Night Versions: The Essential Duran Duran} (1998). The strategy anticipated the remix economy by nearly a decade and maintained the band's dance-floor credibility throughout their pop peak.

\subsubsection{The Nile Rodgers Connection}

Nile Rodgers of Chic --- whose funk bass and rhythm architecture was a primary influence on John Taylor --- was enlisted to remix ``The Reflex'' for its 1984 single release. The remix transformed the album track into the band's first US \#1 hit (also their second UK \#1). This connection, from Chic as influence to Rodgers as collaborator, is one of the cleaner causal chains in Duran Duran's musical history.

\subsubsection{Commercial Records (All-Time)}

\begin{itemize}
  \item Over 100 million records sold worldwide (Billboard / OCC)
  \item 30 top-40 singles in the UK Singles Chart; 14 of them top 10
  \item 21 top-40 singles in the US Billboard Hot 100
  \item Top-five UK albums in five consecutive decades (1980s through 2020s)
  \item US top-ten albums in three decades (1980s, 1990s, 2010s)
  \item 16 studio albums; 50+ singles; 14 video albums; 4 live albums
\end{itemize}

\subsection{Module C Reference: Visual Revolution (1981--1985)}

\subsubsection{The 35mm Innovation}

Prior to Duran Duran's video era, most music videos were shot on video tape with band-playing-on-stage formats. Duran Duran were among the first bands to shoot music videos on 35mm film with professional movie cameras, giving their videos a cinematic polish that was visually indistinguishable from feature film trailers. This decision elevated the perceived production value of the medium itself, and helped transform music videos from promotional afterthoughts into an art form.

\subsubsection{Russell Mulcahy: The Visual Architect}

Australian director Russell Mulcahy directed more than 10 videos for Duran Duran, including ``Hungry Like the Wolf,'' ``Rio,'' ``Save a Prayer,'' ``The Reflex,'' and ``The Wild Boys.'' His stylistic signatures --- fast cuts, tracking shots, glowing lights, neo-noir lighting, windblown drapery --- became the defining visual grammar of MTV's early years. Mulcahy also directed the very first video played on MTV (The Buggles' ``Video Killed the Radio Star,'' 1979), giving his Duran Duran work a particular symbolic resonance.

John Taylor on the collaboration: ``Russell really got us, and we really got him. There was a tremendous kind of mutual addiction that took place.''

The \textit{Rio} video concept was conceived as a ``video album'' --- a revolutionary format at the time. The ``Hungry Like the Wolf,'' ``Save a Prayer,'' and ``Rio'' videos (collectively known as the Sri Lanka trilogy) were shot in Sri Lanka and Antigua, treating the settings as cinematic landscapes rather than backdrops.

\begin{center}
\rowcolors{2}{rowgray}{white}
\begin{tabularx}{\textwidth}{L{3.5cm}L{3cm}XL{2cm}}
\toprule
\rowcolor{primary}
\tableheader{Video} & \tableheader{Location} & \tableheader{Significance} & \tableheader{Year} \\
\midrule
``Girls on Film'' & London studio & Banned by BBC; MTV-edited; first top-10 hit & 1981 \\
``Hungry Like the Wolf'' & Sri Lanka & US breakthrough; MTV heavy rotation & 1982 \\
``Save a Prayer'' & Sri Lanka & Buddhist temples, elephants; defining visual & 1982 \\
``Rio'' & Sri Lanka / Antigua & Yacht sequences; iconic; Grammy for video album & 1982 \\
``Is There Something I Should Know?'' & London soundstage & Deliberate step back: ``Think like The Beatles'' & 1983 \\
``The Reflex'' & Toronto concert & Live footage + Niagara Falls sequence & 1984 \\
``The Wild Boys'' & Shepperton Studios & \textasciitilde\$1M budget; 10-day shoot; Burroughs inspiration & 1984 \\
\bottomrule
\end{tabularx}
\end{center}

\textit{Source: DuranDuran.com, Billboard interview (Mulcahy/J. Taylor), Classic Pop Magazine}

The 1983 Duran Duran Video Album won a Grammy Award for Best Long Form Music Video in 1984.

\subsubsection{Patrick Nagel and the Rio Cover}

For the \textit{Rio} album cover, the band chose American illustrator Patrick Nagel's painting, declining an offer from Andy Warhol (on the grounds that Warhol ``had already done it for The Rolling Stones''). The cover --- a stylized woman's portrait in Nagel's Art Deco-influenced style --- became one of the most iconic album images of the decade. Graphic designer Malcolm Garrett framed the artwork, color-matched typography to the band's attire in interior sleeve designs, and created a unified visual package that extended across singles, posters, and merchandise. The model who inspired Nagel's painting was identified in 2024 as Marcie Hunt, photographed for Vogue Paris (February 1981 edition) by Angelo Tarlazzi.

\subsubsection{Malcolm Garrett: The Graphic Architecture}

Designer Malcolm Garrett worked across the band's early career, creating what \textit{Design Observer} described as a synthesis of Swiss Style modernism and postmodernist pastiche. Key works: the minimalist debut album sleeve; the Nagel-framed \textit{Rio} cover; the constructivist-influenced ``The Reflex'' single; the James Bond-integrated ``A View to a Kill'' typography; and the layered esoteric symbolism of \textit{Seven and the Ragged Tiger} (with Keith Breeden).

\subsection{Module D Reference: Cultural Impact and Legacy (1982--2026)}

\subsubsection{The Second British Invasion}

In July 1983, 20 of the top 40 US singles were by British acts, surpassing the previous record of 14 set in 1965. Rolling Stone devoted a November 1983 issue to the phenomenon with the cover headline ``England Swings: Great Britain invades America's music and style.'' During 1983, 30\% of US record sales were from British acts (Rolling Stone data). Duran Duran occupied the top of the US charts alongside Culture Club, the Police, Kajagoogoo, and Eddy Grant.

Parke Puterbaugh wrote in Rolling Stone (November 1983): ``Next to the prosaic, foursquare appearance of the American bands, such acts as Duran Duran seemed like caviar.''

\subsubsection{Beatlemania-Scale Fame (1983--1984)}

By 1984, the band's fame had escalated to levels their members compared to Beatlemania. Public appearances produced screaming mobs. Tokyo arena shows sold out within minutes of announcement. Princess Diana declared Duran Duran her favourite band. People magazine called them ``the prettiest boys in rock.'' The British press dubbed them ``the Fab Five.''

Simon Le Bon's yacht Drum capsized during the 1985 Fastnet race; Le Bon and crew members were trapped underwater for 40 minutes before rescue --- an incident that illustrated the band's unusual position in mainstream celebrity culture.

\subsubsection{Critical Reception and the Bee Gees Curse}

Duran Duran suffered what Moby described as ``the Bee Gees curse'': write amazing songs, sell enormous quantities of records, and incur critical dismissal from the rock establishment. During their 1980s peak, critics frequently categorized them as manufactured pop product. Sunday Herald noted: ``To describe them, as some have, as the first Boy Band misrepresents their appeal. Their weapons were never just their looks, but self-penned songs.'' Critical reappraisal has continued from the 1990s onward, with \textit{Rio} receiving a 90/100 critic score on Album of the Year and recognition as a defining document of the era.

\subsubsection{Awards and Recognition}

\begin{itemize}
  \item Two Brit Awards (including 2004 Outstanding Contribution to Music)
  \item Two Grammy Awards (Best Long Form Music Video 1984; second Grammy at other ceremony)
  \item MTV Video Music Award for Lifetime Achievement
  \item MTV Europe Music Awards Video Visionary Award
  \item Hollywood Walk of Fame star
  \item Rock and Roll Hall of Fame induction, 2022 (class of 2022)
\end{itemize}

\subsubsection{The 2022 Rock and Roll Hall of Fame Induction}

Inducted in 2022, the induction was marked by Andy Taylor's absence due to advanced prostate cancer --- a diagnosis he revealed in a letter read at the ceremony. Taylor had missed the induction; the remaining members (Le Bon, Rhodes, John Taylor, Roger Taylor) accepted. The moment foregrounded the band's resilience across decades of personnel challenges and health crises.

\subsubsection{Documented Influence on Later Artists}

\begin{itemize}
  \item \textbf{The Killers} --- cited Duran Duran as direct influence on their synthesizer-driven rock approach
  \item \textbf{Franz Ferdinand} --- acknowledged the Duran Duran template for dance-floor-oriented pop rock
  \item \textbf{The Bravery} --- cited the band's visual-first identity as foundational
  \item \textbf{Moby} --- website diary (2003): ``They were cursed by... the Bee Gees curse''
  \item \textbf{Nick Rhodes} (2021): ``You were proud of the fact that you didn't sound exactly like somebody else'' --- describing era alongside The Cure, INXS, Depeche Mode, all maintaining distinct identities
\end{itemize}

\textit{Sources: Rate Your Music compilation; Saturday Evening Post retrospective; Classic Pop Magazine}

\subsection{Safety and Sensitivity Considerations}

\begin{center}
\rowcolors{2}{rowgray}{white}
\begin{tabularx}{\textwidth}{L{4cm}XC{2.5cm}}
\toprule
\rowcolor{primary}
\tableheader{Consideration} & \tableheader{Guidance} & \tableheader{Type} \\
\midrule
Andy Taylor's cancer & Handle with care; confirm current status before publication & GATE \\
``Girls on Film'' video content & Factual documentation is appropriate; avoid editorial judgment on 1981 standards applied retroactively & ANNOTATE \\
Chart records & All claims require specific source citation (Billboard, OCC) & GATE \\
Critical dismissal framing & Present evidence fairly; avoid advocacy for or against critical positions & ANNOTATE \\
Sales figures & Multiple sources cite 70M and 100M at different points; note the timeline (2016 vs later) & GATE \\
\bottomrule
\end{tabularx}
\end{center}

\subsection{Source Bibliography}

\subsubsection{Primary and Official Sources}
\begin{itemize}
  \item DuranDuran.com --- Official Band Timeline 1980--1995 (accessed 2026)
  \item DuranDuran.com --- Full Timeline (accessed 2026)
  \item Rock and Roll Hall of Fame --- Class of 2022 inductee documentation
  \item RIAA --- Certification database (platinum/gold records)
\end{itemize}

\subsubsection{Music Industry Publications}
\begin{itemize}
  \item Billboard --- Chart data and ``over 100 million records sold'' commercial record
  \item Rolling Stone --- Parke Puterbaugh, ``Anglomania: The Second British Invasion,'' November 1983
  \item Classic Pop Magazine --- ``Album By Album: Duran Duran'' (2023 retrospective)
  \item Classic Pop Magazine --- ``Making Duran Duran: Rio'' (2021)
  \item Saturday Evening Post --- ``What Makes Duran Duran So Durable?'' (2021)
  \item This Day in Music --- Duran Duran artist entry
\end{itemize}

\subsubsection{Reference and Encyclopedic Sources}
\begin{itemize}
  \item Wikipedia --- ``Duran Duran'' (cross-referenced, March 2026)
  \item Wikipedia --- ``Duran Duran discography'' (March 2026)
  \item Wikipedia --- ``Second British Invasion'' (cross-referenced)
  \item Wikipedia --- ``New Romantic'' (cross-referenced)
  \item Wikipedia --- ``Russell Mulcahy'' (March 2026)
  \item Last.fm --- Duran Duran artist biography
  \item Rate Your Music --- Duran Duran artist page and critical consensus
  \item Album of the Year --- Duran Duran discography scores
\end{itemize}

\newpage

%% ============================================================
\stageheader{STAGE 3 --- MISSION PACKAGE}{tertiary}

\section{Milestone Specification}

\subsection{Mission Objective}

Produce a comprehensive, publication-quality research document on Duran Duran covering four modules: Origins and Formation (1978--1981), Musical Architecture (1981--1993), Visual Revolution (1981--1985), and Cultural Impact and Legacy (1982--2026). The final deliverable will serve as a definitive reference for GSD ecosystem contributors, providing sourced, structured knowledge across the band's full history and significance.

\subsection{Architecture Overview}

\begin{asciibox}
WAVE EXECUTION ARCHITECTURE
======================================

WAVE 0 [FOUNDATION]
  |-- F-01: Shared schemas and data models
  |-- F-02: Source index and citation registry
  |-- F-03: Output templates
  Sequential | Haiku | ~0.5 context windows

WAVE 1A [ORIGINS + MUSIC]            WAVE 1B [VISUAL + LEGACY]
  |-- MOD-A: Origins (Sonnet)          |-- MOD-C: Visual Rev (Opus)
  |-- MOD-B: Music Arch (Sonnet)       |-- MOD-D: Legacy (Opus)
  Parallel | ~2.0 context windows      Parallel | ~2.0 context windows

WAVE 2 [SYNTHESIS]
  |-- SYN-01: Cross-module integration (Opus)
  |-- SYN-02: Through-line weaving
  Sequential | ~1.0 context window

WAVE 3 [PUBLICATION + VERIFICATION]
  |-- PUB-01: Final document assembly (Sonnet)
  |-- VER-01: Source validation (Sonnet)
  |-- VER-02: Accuracy check (Sonnet)
  Sequential | ~1.0 context window
\end{asciibox}

\subsection{System Layers}

\begin{enumerate}
  \item \textbf{Foundation Layer} --- Shared data schemas, citation registry, output format templates
  \item \textbf{Research Layer} --- Four parallel module surveys (Origins, Music, Visual, Legacy)
  \item \textbf{Synthesis Layer} --- Cross-module integration, relationship mapping, narrative weaving
  \item \textbf{Publication Layer} --- Final document assembly, source validation, quality verification
\end{enumerate}

\subsection{Deliverables}

\begin{center}
\rowcolors{2}{rowgray}{white}
\begin{tabularx}{\textwidth}{C{0.8cm}L{4cm}XL{3cm}}
\toprule
\rowcolor{primary}
\tableheader{\#} & \tableheader{Deliverable} & \tableheader{Acceptance Criteria} & \tableheader{Spec} \\
\midrule
D-01 & Origins and Formation Survey & All 12 lineup phases documented; Rum Runner fully traced & MOD-A complete \\
D-02 & Musical Architecture Analysis & 16 albums catalogued; sound system analyzed & MOD-B complete \\
D-03 & Visual Revolution Study & 7+ videos detailed; Mulcahy/Nagel/Garrett documented & MOD-C complete \\
D-04 & Cultural Legacy Assessment & Second British Invasion data; Hall of Fame; 3+ influence citations & MOD-D complete \\
D-05 & Integrated Research Document & All four modules woven; cross-references resolved & SYN-01 output \\
D-06 & Verified Source Register & All claims attributed; no unsourced statistics & VER-01 output \\
\bottomrule
\end{tabularx}
\end{center}

\subsection{Component Breakdown}

\begin{center}
\rowcolors{2}{rowgray}{white}
\begin{tabularx}{\textwidth}{L{3cm}L{3cm}C{2cm}C{2cm}}
\toprule
\rowcolor{primary}
\tableheader{Component} & \tableheader{Dependencies} & \tableheader{Model} & \tableheader{Est. Tokens} \\
\midrule
F-01 Schema Setup & None & Haiku & 2K \\
F-02 Source Index & F-01 & Haiku & 3K \\
MOD-A Origins Survey & F-01, F-02 & Sonnet & 25K \\
MOD-B Music Architecture & F-01, F-02 & Sonnet & 35K \\
MOD-C Visual Revolution & F-01, F-02 & Opus & 30K \\
MOD-D Legacy Assessment & F-01, F-02 & Opus & 30K \\
SYN-01 Integration & MOD-A through D & Opus & 20K \\
PUB-01 Assembly & SYN-01 & Sonnet & 15K \\
VER-01 Source Validation & PUB-01 & Sonnet & 10K \\
\bottomrule
\end{tabularx}
\end{center}

\subsection{Model Assignment Rationale}

\begin{center}
\rowcolors{2}{rowgray}{white}
\begin{tabularx}{\textwidth}{L{2cm}C{1.5cm}XL{3cm}}
\toprule
\rowcolor{primary}
\tableheader{Model} & \tableheader{\%} & \tableheader{Rationale} & \tableheader{Assigned To} \\
\midrule
Opus & 30\% & Cultural analysis, cross-module synthesis, legacy judgment, critical framing & MOD-C, MOD-D, SYN-01 \\
Sonnet & 60\% & Survey compilation, discography documentation, structured cataloguing & MOD-A, MOD-B, PUB-01, VER-01 \\
Haiku & 10\% & Schema setup, template generation, scaffolding & F-01, F-02 \\
\bottomrule
\end{tabularx}
\end{center}

\section{Wave Execution Plan}

\subsection{Wave Summary}

\begin{center}
\rowcolors{2}{rowgray}{white}
\begin{tabularx}{\textwidth}{C{1.5cm}L{3cm}C{2cm}C{2cm}C{1.5cm}X}
\toprule
\rowcolor{primary}
\tableheader{Wave} & \tableheader{Name} & \tableheader{Mode} & \tableheader{Model} & \tableheader{Windows} & \tableheader{Output} \\
\midrule
Wave 0 & Foundation & Sequential & Haiku & 0.5 & Schemas, index, templates \\
Wave 1A & Origins + Music & Parallel & Sonnet & 2.0 & MOD-A, MOD-B surveys \\
Wave 1B & Visual + Legacy & Parallel & Opus & 2.0 & MOD-C, MOD-D surveys \\
Wave 2 & Synthesis & Sequential & Opus & 1.0 & Integrated document draft \\
Wave 3 & Publication & Sequential & Sonnet & 1.0 & Final verified document \\
\bottomrule
\end{tabularx}
\end{center}

\subsection{Wave 0: Foundation}

\begin{center}
\rowcolors{2}{rowgray}{white}
\begin{tabularx}{\textwidth}{L{2cm}L{3.5cm}XC{2cm}}
\toprule
\rowcolor{primary}
\tableheader{Task ID} & \tableheader{Task} & \tableheader{Description} & \tableheader{Model} \\
\midrule
F-01 & Schema Setup & Define shared data types: AlbumRecord, SingleRecord, VideoRecord, PersonRecord, AwardRecord & Haiku \\
F-02 & Source Index & Build citation registry from Stage 2 bibliography; assign citation keys & Haiku \\
F-03 & Templates & Generate output templates for all four module survey formats & Haiku \\
\bottomrule
\end{tabularx}
\end{center}

\subsection{Wave 1A: Origins + Music Architecture (Parallel)}

\begin{center}
\rowcolors{2}{rowgray}{white}
\begin{tabularx}{\textwidth}{L{2cm}L{3.5cm}XC{2cm}}
\toprule
\rowcolor{primary}
\tableheader{Task ID} & \tableheader{Task} & \tableheader{Description} & \tableheader{Model} \\
\midrule
A-01 & Lineup Evolution & Document all lineup phases 1978--1980 with sources & Sonnet \\
A-02 & Rum Runner Analysis & Trace Berrow brothers, management structure, lab function & Sonnet \\
A-03 & EMI Signing Context & Bidding war, New Romantic positioning, debut recording & Sonnet \\
A-04 & Debut Album Deep Dive & Track-by-track analysis; influences; reception & Sonnet \\
B-01 & Sound System Map & Member roles, influences, interaction model & Sonnet \\
B-02 & Discography Survey & All 16 albums: chart data, key singles, production notes & Sonnet \\
B-03 & Night Versions Study & Remixing philosophy; 12'' culture; dance-floor strategy & Sonnet \\
B-04 & Nile Rodgers Thread & Chic influence to ``Reflex'' remix; the Rodgers connection & Sonnet \\
\bottomrule
\end{tabularx}
\end{center}

\subsection{Wave 1B: Visual Revolution + Cultural Legacy (Parallel)}

\begin{center}
\rowcolors{2}{rowgray}{white}
\begin{tabularx}{\textwidth}{L{2cm}L{3.5cm}XC{2cm}}
\toprule
\rowcolor{primary}
\tableheader{Task ID} & \tableheader{Task} & \tableheader{Description} & \tableheader{Model} \\
\midrule
C-01 & 35mm Innovation & Document the technical shift; compare pre/post DD video aesthetics & Opus \\
C-02 & Mulcahy Profile & Full collaboration history; 10+ videos; production methods & Opus \\
C-03 & Sri Lanka Trilogy & Detailed analysis: locations, narrative, cultural significance & Opus \\
C-04 & Nagel/Garrett Visual System & Rio cover, sleeve design philosophy, unified identity & Opus \\
D-01 & Second British Invasion & Quantitative data; chart records; Rolling Stone coverage & Opus \\
D-02 & Beatlemania Era & 1983--84 cultural fever; Princess Diana; press documentation & Opus \\
D-03 & Longevity Analysis & Five decades of top-five UK albums; resilience mechanisms & Opus \\
D-04 & Hall of Fame + Andy Taylor & 2022 induction; cancer revelation; emotional context & Opus \\
D-05 & Influence Documentation & The Killers, Franz Ferdinand, Bravery, Moby citations & Opus \\
\bottomrule
\end{tabularx}
\end{center}

\subsection{Wave 2: Synthesis}

\begin{center}
\rowcolors{2}{rowgray}{white}
\begin{tabularx}{\textwidth}{L{2cm}L{3.5cm}XC{2cm}}
\toprule
\rowcolor{primary}
\tableheader{Task ID} & \tableheader{Task} & \tableheader{Description} & \tableheader{Model} \\
\midrule
SYN-01 & Cross-Module Integration & Weave all four modules; resolve contradictions; trace connections & Opus \\
SYN-02 & Narrative Arc & Through-line: Birmingham art school $\rightarrow$ MTV $\rightarrow$ Hall of Fame & Opus \\
SYN-03 & Critical Framing & The Bee Gees Curse; substance beneath the style argument & Opus \\
\bottomrule
\end{tabularx}
\end{center}

\subsection{Wave 3: Publication and Verification}

\begin{center}
\rowcolors{2}{rowgray}{white}
\begin{tabularx}{\textwidth}{L{2cm}L{3.5cm}XC{2cm}}
\toprule
\rowcolor{primary}
\tableheader{Task ID} & \tableheader{Task} & \tableheader{Description} & \tableheader{Model} \\
\midrule
PUB-01 & Document Assembly & Final integrated document in target format & Sonnet \\
VER-01 & Source Validation & SC-SRC: all citations traceable to professional sources & Sonnet \\
VER-02 & Numerical Attribution & SC-NUM: all chart positions and sales figures attributed & Sonnet \\
VER-03 & Consistency Check & Cross-module data consistency; no contradictions & Sonnet \\
\bottomrule
\end{tabularx}
\end{center}

\subsection{Cache Optimization Strategy}

\begin{itemize}
  \item Wave 0 schema outputs should be cached at session start and referenced by all Wave 1 components
  \item The Stage 2 Research Reference (this document) should be loaded once and cached for the full Wave 1 execution window
  \item Wave 1A and 1B can run in separate sessions with the shared schemas as the cache anchor
  \item Opus synthesis in Wave 2 benefits most from a single long context window containing all four module outputs
  \item Target cache TTL exploitation: load all Wave 1 outputs into Wave 2's initial context in a single burst
\end{itemize}

\subsection{Token Budget}

\begin{center}
\rowcolors{2}{rowgray}{white}
\begin{tabularx}{\textwidth}{L{4cm}C{2cm}C{2cm}C{2cm}}
\toprule
\rowcolor{primary}
\tableheader{Wave} & \tableheader{Opus (K)} & \tableheader{Sonnet (K)} & \tableheader{Haiku (K)} \\
\midrule
Wave 0 & 0 & 0 & 5 \\
Wave 1A & 0 & 60 & 0 \\
Wave 1B & 60 & 0 & 0 \\
Wave 2 & 20 & 0 & 0 \\
Wave 3 & 0 & 25 & 0 \\
\midrule
\textbf{Total} & \textbf{80K} & \textbf{85K} & \textbf{5K} \\
\textbf{Percentage} & \textbf{47\%} & \textbf{50\%} & \textbf{3\%} \\
\bottomrule
\end{tabularx}
\end{center}

\textit{Note: Opus usage is slightly elevated above the 30\% guideline due to the cultural analysis and synthesis demands of Modules C and D. This is appropriate given the depth of critical framing required.}

\subsection{Dependency Graph}

\begin{asciibox}
DEPENDENCY GRAPH
======================================

F-01 ----+----> MOD-A (A-01..A-04)
F-02 ----|        |
F-03 ----+----> MOD-B (B-01..B-04)  ----+
         |                               |
         +----> MOD-C (C-01..C-04)       |
         |                               |---> SYN-01 ---> PUB-01 ---> D-06 (Done)
         +----> MOD-D (D-01..D-05) ----+

PARALLEL GATES:
  [Wave 0 complete] -> [Wave 1A + Wave 1B can start simultaneously]
  [Wave 1A complete] -> [Wave 2 can start when Wave 1B also complete]
  [Wave 2 complete] -> [Wave 3 begins]
\end{asciibox}

\section{Test Plan}

\subsection{Test Categories}

\begin{center}
\rowcolors{2}{rowgray}{white}
\begin{tabularx}{\textwidth}{L{3.5cm}XC{2cm}C{2.5cm}C{1.5cm}}
\toprule
\rowcolor{primary}
\tableheader{Category} & \tableheader{Purpose} & \tableheader{Count} & \tableheader{Failure Action} & \tableheader{\%} \\
\midrule
Safety-Critical & Non-negotiable source and accuracy gates & 5 & BLOCK & 15\% \\
Core Functionality & Deliverable acceptance per success criteria & 18 & BLOCK & 45\% \\
Integration & Cross-module validation & 8 & BLOCK & 25\% \\
Edge Cases & Robustness and completeness & 5 & LOG & 15\% \\
\bottomrule
\end{tabularx}
\end{center}

\subsection{Safety-Critical Tests}

\begin{center}
\rowcolors{2}{rowgray}{white}
\begin{tabularx}{\textwidth}{C{1.5cm}L{4.5cm}X}
\toprule
\rowcolor{primary}
\tableheader{Test ID} & \tableheader{Verifies} & \tableheader{Expected Behavior} \\
\midrule
SC-SRC & Source quality gate & All citations traceable to professional music publications, official band sources, or established encyclopedic references. Zero entertainment blogs or unsourced fan sites. \\
SC-NUM & Numerical attribution & Every chart position (UK, US), sales figure (100M records), and percentage claim attributed to a specific named source (Billboard, OCC, RIAA, Rolling Stone). \\
SC-ADV & No advocacy & Document presents critical reception evidence (Bee Gees Curse, critical dismissal) without advocating for a particular critical position. Evidence presented; judgment left to reader. \\
SC-MED & Andy Taylor health information & Andy Taylor's prostate cancer disclosure handled accurately; current health status verified before publication; no speculation beyond confirmed reports. \\
SC-CON & Content sensitivity & ``Girls on Film'' video documented factually with historical context; no anachronistic moralizing applied to 1981 production standards. \\
\bottomrule
\end{tabularx}
\end{center}

\subsection{Core Functionality Tests (Selected)}

\begin{center}
\rowcolors{2}{rowgray}{white}
\begin{tabularx}{\textwidth}{C{1.5cm}L{4.5cm}X}
\toprule
\rowcolor{primary}
\tableheader{Test ID} & \tableheader{Verifies} & \tableheader{Expected Behavior} \\
\midrule
CF-01 & Origins lineup completeness & All pre-classic lineup phases documented: Duffy era, Wickett era, Thomas era, with dates \\
CF-02 & Rum Runner structure & Document identifies Berrows, Tritec Music, specific roles of each band member at the club \\
CF-03 & Discography completeness & All 16 studio albums present with UK/US chart positions and at least one key single noted \\
CF-04 & Sound system roles & Each of the five classic members has their specific musical role, influences, and sonic contribution documented \\
CF-05 & Video detail depth & At least 7 videos documented with location, year, director, and significance \\
CF-06 & Mulcahy collaboration & More than 10 video titles credited to Mulcahy; production approach characterized \\
CF-07 & Patrick Nagel documentation & Rio cover origin documented; Andy Warhol declined story included; 2024 model identification noted \\
CF-08 & Malcolm Garrett system & At least 4 Garrett-designed covers described; design philosophy characterized \\
CF-09 & Second British Invasion data & July 1983 chart data included; 30\% US sales figure cited from Rolling Stone \\
CF-10 & Beatlemania framing & Princess Diana reference; People magazine quote; member comparisons to Beatles documented \\
CF-11 & Commercial records & 100M records, 30 UK top-40 singles, 21 US top-40 singles all present with Billboard/OCC attribution \\
CF-12 & Awards complete & 2 Brits, 2 Grammys, MTV Lifetime Achievement, Hollywood Walk of Fame, Hall of Fame all documented \\
CF-13 & Hall of Fame details & 2022 class; Andy Taylor absence; cancer letter at ceremony documented \\
CF-14 & Influence citations & The Killers, Franz Ferdinand, The Bravery all cited with specific quote or attribution \\
CF-15 & Night Versions strategy & Remixing philosophy documented; 12'' club culture; Night Versions compilation cited \\
CF-16 & Nile Rodgers thread & Chic as influence; ``The Reflex'' remix; US \#1 outcome documented \\
CF-17 & Five decades claim & Top-five UK albums in five consecutive decades verified with specific album examples \\
CF-18 & Wedding Album comeback & 1993 self-titled album (``Wedding Album'') documented; ``Ordinary World'' and ``Come Undone'' chart positions noted \\
\bottomrule
\end{tabularx}
\end{center}

\subsection{Integration Tests}

\begin{center}
\rowcolors{2}{rowgray}{white}
\begin{tabularx}{\textwidth}{C{1.5cm}L{4.5cm}X}
\toprule
\rowcolor{primary}
\tableheader{Test ID} & \tableheader{Verifies} & \tableheader{Expected Behavior} \\
\midrule
IN-01 & A $\rightarrow$ C cascade & Document traces the art school origins (Module A) to visual priority (Module C); the Rum Runner aesthetic feeds into video ambition \\
IN-02 & B $\rightarrow$ C cascade & Night Versions remixing strategy (Module B) connects to MTV rotation strategy (Module C); club sound precedes video strategy \\
IN-03 & C $\rightarrow$ D cascade & Music video revolution (Module C) explicitly connects to Second British Invasion success (Module D) \\
IN-04 & B $\rightarrow$ D longevity & Wedding Album comeback (Module B) connects to longevity paradox resolution (Module D) \\
IN-05 & Cross-module consistency & Sales figures, chart positions, and award lists are consistent across all four modules; no contradictions \\
IN-06 & Name consistency & ``Andy Taylor,'' ``Roger Taylor,'' ``John Taylor'' consistently differentiated across all modules; no confusion of roles \\
IN-07 & Timeline consistency & Formation year (1978), classic lineup (1980), US breakthrough (1983), lineup split (1986), reunion (2001), Hall of Fame (2022) consistent across all modules \\
IN-08 & Self-containment & A reader with no prior knowledge of Duran Duran can read the document and understand the band's full arc; no undefined terms or unexplained references \\
\bottomrule
\end{tabularx}
\end{center}

\subsection{Verification Matrix}

\begin{center}
\rowcolors{2}{rowgray}{white}
\begin{tabularx}{\textwidth}{XL{3.5cm}C{1.5cm}}
\toprule
\rowcolor{primary}
\tableheader{Success Criterion} & \tableheader{Test ID(s)} & \tableheader{Status} \\
\midrule
1. Origins module documents all lineup phases with sources & CF-01, IN-07 & Pending \\
2. Rum Runner's role fully traced with specifics & CF-02, IN-01 & Pending \\
3. Full discography covers all 16 albums & CF-03, CF-17 & Pending \\
4. Sound architecture identifies roles of each member & CF-04, IN-06 & Pending \\
5. Visual Revolution documents Mulcahy for 5+ videos & CF-05, CF-06 & Pending \\
6. Nagel and Garrett analyzed as unified visual identity & CF-07, CF-08 & Pending \\
7. MTV strategy documented with chart correlation data & CF-09, IN-03 & Pending \\
8. Second British Invasion with quantitative data & CF-09, CF-10 & Pending \\
9. Legacy cites 3+ named contemporary bands & CF-14 & Pending \\
10. All chart claims attributed to specific sources & SC-NUM, CF-11 & Pending \\
11. Lineup turbulence documented with dates and consequences & CF-01, IN-07 & Pending \\
12. Andy Taylor's 2022 absence documented & CF-13, SC-MED & Pending \\
\bottomrule
\end{tabularx}
\end{center}

\section{Mission Crew Manifest}

\begin{center}
\rowcolors{2}{rowgray}{white}
\begin{tabularx}{\textwidth}{L{2cm}L{3cm}X}
\toprule
\rowcolor{primary}
\tableheader{Role} & \tableheader{Agent} & \tableheader{Responsibility} \\
\midrule
FLIGHT & Orchestrator & Overall mission direction; wave go/no-go gates; human escalation \\
PLAN & Planner & Wave decomposition; parallel track optimization; dependency management \\
SCOUT & Research Agent & Additional web research for gap-filling; source discovery \\
EXEC-1 & Executor Alpha & Wave 1A document production (MOD-A, MOD-B) \\
EXEC-2 & Executor Beta & Wave 1B document production (MOD-C, MOD-D) \\
CRAFT-MUS & Music Specialist & Musical architecture depth: theory, production, discography \\
CRAFT-HIST & History Specialist & Cultural context: Birmingham scene, New Romantic, Second British Invasion \\
VERIFY & Verification Agent & Source validation; accuracy checking; SC test enforcement \\
INTEG & Integration Agent & Cross-module relationship validation; consistency checks \\
CAPCOM & HITL Interface & Human approval at wave boundaries; escalation channel \\
BUDGET & Resource Manager & Token budget tracking; session planning; model allocation \\
LOG & Chronicle Agent & Audit trail; commit messages; milestone summaries \\
\bottomrule
\end{tabularx}
\end{center}

\section{Execution Summary}

\begin{center}
\rowcolors{2}{rowgray}{white}
\begin{tabularx}{\textwidth}{L{5cm}X}
\toprule
\rowcolor{primary}
\tableheader{Metric} & \tableheader{Value} \\
\midrule
Topic & Duran Duran: Sound, Vision, and the Architecture of Pop \\
Activation Profile & Squadron (12 roles) \\
Research Modules & 4 (Origins, Music, Visual, Legacy) \\
Total Components & 9 (2 foundation + 4 module + 1 synthesis + 2 publication) \\
Estimated Context Windows & 6--8 across 3 sessions \\
Total Test Count & 36 (5 SC + 18 CF + 8 IN + 5 EC) \\
Model Split & Opus 47\% / Sonnet 50\% / Haiku 3\% \\
Wave Count & 5 waves (0 through 3, with 1A/1B parallel) \\
Estimated Token Budget & 170K total (80K Opus / 85K Sonnet / 5K Haiku) \\
Human Approval Gates & Wave 0 complete; Wave 1 complete; Wave 2 complete \\
Status & Ready for GSD Orchestrator handoff \\
\bottomrule
\end{tabularx}
\end{center}

\vspace{2em}

\begin{quotebox}
{\large\sffamily\textcolor{primary}{``There wasn't a day that passed when we didn't try to move ourselves forward a little bit --- and that's what paid off.''}}

\vspace{0.5em}
{\small\sffamily\textcolor{subtlegray}{--- Nick Rhodes, on making Rio (The Guardian interview)}}

\vspace{1em}
{\small\itshape\textcolor{darktext}{The Amiga Principle, stated in plain English: incremental, principled building produces staggering outcomes. Five members from Birmingham, a nightclub, a director who shot the first MTV video, a painter named Nagel, and a cable channel that nobody yet knew would reshape the world. None of those nodes was exceptional in isolation. The architecture was everything. The space between them was where the music lived.}}
\end{quotebox}

\end{document}
